\par\Needspace{4\baselineskip}
\begin{center}
\hypertarget{native-042-article025}{}\pdfbookmark[2]{25.}{article-025}25.
\end{center}
To solve the inverse problem of determining the true anomaly and the radius vector from the time, we start from $N=\lambda kb^{-\frac{3}{2}}t$ and use equation XI to obtain the auxiliary quantity $u$ or $F$. This transcendental equation must be solved by trial, and the work can be shortened by methods analogous to those we explained in Article 11. We shall dispense with a fuller explanation. Hyperbolic motion may occur only very rarely in the heavens, so it seems hardly worth refining the rules for it as carefully as those for elliptic motion. \SourceLeaf{043}Besides, every case that might arise can be dealt with by another method to be given below. Once $F$ or $u$ has been found, $v$ is determined from it by formula III, and then $r$ by either II or VIII. It is even more convenient to obtain $v$ and $r$ together by formulas VI and VII. If desired, one or another of the remaining formulas can be used to check the calculation.

\par\Needspace{4\baselineskip}
\begin{center}
\hypertarget{native-043-article026}{}\pdfbookmark[2]{26.}{article-026}26.
\end{center}
Example. With $e$ and $b$ as in the preceding example, let $t=65{,}41236$. We seek $v$ and $r$. Using \textsc{Briggs}ian logarithms, we have
\begin{center}
\begin{tabular}{@{}ll@{}}
\makebox[18em][l]{$\log t$\hfill .\,.\,.\,.\,.\,.\,.\,.\,.\,.\,.\,.} & $1{,}815\,6598$\\[.3em]
\makebox[18em][l]{$\log\lambda kb^{-\frac{3}{2}}$\hfill .\,.\,.\,.\,.\,.\,.\,.} & $6{,}970\,2758$\\[.3em]
\hline
\makebox[18em][l]{$\log N$\hfill .\,.\,.\,.\,.\,.\,.\,.\,.\,.\,.} & $8{,}785\,9356$,
\end{tabular}\\[.25em]
whence $N=0{,}061\,0851\,4$.
\end{center}
We thus find that the equation $N=\lambda e\operatorname{tang}F-\log\operatorname{tang}(45^0+\frac{1}{2}F)$ is satisfied by
\[
F=25^0\,24^{\prime}\,27^{\prime\prime}{,}66,
\]
from which formula III gives
\begin{center}
\begin{tabular}{@{}ll@{}}
\makebox[18em][l]{$\log\operatorname{tang}\frac{1}{2}F$\hfill .\,.\,.\,.\,.\,.\,.} & $9{,}353\,0120$\\[.3em]
\makebox[18em][l]{$\log\operatorname{tang}\frac{1}{2}\psi$\hfill .\,.\,.\,.\,.\,.\,.} & $9{,}531\,8179$\\[.3em]
\hline
\makebox[18em][l]{$\log\operatorname{tang}\frac{1}{2}v$\hfill .\,.\,.\,.\,.\,.\,.} & $9{,}821\,1941$
\end{tabular}
\end{center}
and hence $\frac{1}{2}v=33^0\,31^{\prime}\,29^{\prime\prime}{,}89$ and $v=67^0\,2^{\prime}\,59^{\prime\prime}{,}78$. We further obtain
\begin{center}
\footnotesize
\begin{tabular}{@{}ll@{}}
$\left.\begin{array}{ll}
\text{C.}\log\cos\frac{1}{2}(v+\psi)\quad .\,.\,. &0{,}213\,7476\\[.3em]
\text{C.}\log\cos\frac{1}{2}(v-\psi)\quad .\,.\,. &0{,}014\,5197
\end{array}\right\}$ & \begin{tabular}{@{}ll@{}}
difference .\,.\,.\,.\,.\,.\,.\,.\,.\,.\,.\,. & $0{,}199\,2279$\\[.3em]
$\log\operatorname{tang}(45^0+\frac{1}{2}F)$ .\,.\,. & $0{,}199\,2280$
\end{tabular}
\end{tabular}
\end{center}
\begin{center}
\begin{tabular}{@{}ll@{}}
\makebox[18em][l]{$\log\frac{p}{2e}$\hfill .\,.\,.\,.\,.\,.\,.\,.\,.\,.\,.} & $9{,}972\,5868$\\[.3em]
\hline
\makebox[18em][l]{$\log r$\hfill .\,.\,.\,.\,.\,.\,.\,.\,.\,.\,.\,.} & $0{,}200\,8541$
\end{tabular}
\end{center}
